% 5. Discussion
\section{Discussion}

\paragraph{Main finding 1: the GABA-specific hypothesis is not supported.}
The pre-registered test --- degree matching, degree-controlled regression,
and an ensemble of 100 degree-preserving null networks --- did not provide
evidence that GABAergic identity independently predicts structural control
impact beyond degree-related structural effects. The small matched-pair
advantage ($\delta\approx0.10$--$0.11$ depending on panel size) lies
inside the null ensemble, and the regression residual is indistinguishable
from zero. This is a meaningful negative result: given the documented
GABA--degree association \cite{lin2024}, the analysis asked precisely
whether anything remains beyond degree, and at this resolution and
estimator, nothing does. The claim is bounded: it does not establish that
degree \emph{causes} control impact, nor that neurotransmitters play no
role in circuit function beyond what a structural, shortest-path analysis
can detect.

\paragraph{Main finding 2: structural chokepoints concentrate in
visual-centrifugal architecture.}
The strongest neurons are not scattered randomly across the degree
distribution: they concentrate in visual-centrifugal circuitry
(2.63$\times$ after degree matching, $z=5.3$), consistently across
threshold choices ($K=25/50/100$). This is the study's positive
contribution --- a concrete, ranked, anatomy-linked structural chokepoint
map of the adult fly brain, quantified against degree-matched nulls rather
than raw base rates. It is consistent with anatomical knowledge that
centrifugal feedback neurons bridge the massive optic-lobe periphery to
the central brain \cite{matsliah2024,hoeller2026,mickels2026}: such bridge
positions make single-neuron removal disproportionately damaging to
eye--brain communication paths.

\paragraph{Main finding 3: some individuals exceed degree-matched peers
substantially.}
Thirteen of the top 50 beat their degree-matched peers individually, and
one moderate-degree centrifugal neuron reaches ${\sim}156\times$ its peer
median. Degree alone therefore does not completely explain the observed
architecture. As clearly labeled hypotheses for future work, plausible
contributors include inter-batch (inter-neuropil) bridging topology,
bilateral mirror-position redundancy, and cell-type-specific wiring
patterns; none of these was tested here, and we do not assert any as
established mechanisms.

\paragraph{Biological meaning.}
All findings concern structural network features: sensitivity of a
shortest-path efficiency functional to computational node removal. They
identify \emph{network-control-sensitive neurons}, not functional
necessity; removal impact is an abstraction of silencing, not a
physiological manipulation. The results generate specific, testable
predictions --- e.g., that perturbing the identified centrifugal
chokepoints should disproportionately affect eye--brain signal
propagation --- but establishing that requires functional experiments.

\paragraph{Limitations.}
(1) Structural, not functional: no firing rates, dynamics, or state
dependence. (2) Dependence on a single reconstruction (FAFB v783, one
female brain); cross-dataset replication (BANC, MaleCNS) remains future
work. (3) The node-removal metric is one abstraction; a reachability-drop
metric correlates only moderately ($\rho=0.39$) and both are reported.
(4) Neurotransmitter annotations are predictions, not measured signs;
GLUT was deliberately analyzed separately. (5) Degree matching is greedy
1:1 within $\pm10\%$; residual degree imbalance within windows cannot be
fully excluded, though the null ensemble --- which preserves degree
exactly --- addresses the dominant confound. (6) The full no-threshold
connection-table sensitivity analysis was RAM-infeasible here
(${\sim}50$M pairs; deferred to high-RAM), and the Buhmann table is
non-comparable; this axis remains open. (7) Empirical $p$-values on 100
nulls have ${\sim}0.01$ tail resolution, and the null comparison uses the
$k{=}4$ panel (observed $\delta=0.098$) while the primary matched-pair
statistic reads $\delta=0.111$ at $k{=}8$; both statistics agree on the
verdict, and the OLS independently points the same way.

\paragraph{Future work.}
Functional recordings targeting the identified chokepoints; activity-aware
and temporal-dynamics models; synaptic-weight-aware efficiency metrics;
higher-resolution (no-threshold) sensitivity analyses on high-RAM
hardware; cross-dataset replication; and experimental validation of the
centrifugal chokepoint predictions.
